\clearpage
\oldpage[II]{217}
\markboth{ERRATA ET ADDENDA, LISTE 1}{ERRATA ET ADDENDA, LISTE 1}
\phantomsection
\addcontentsline{toc}{section}{Errata et Addenda (Liste 1)}
\section*{Errata et Addenda}
\thispagestyle{plain}
\begin{center}
\large (Liste 1)
\end{center}

\subsection*{Chapitre 0}

\paragraph{(0, 2.1.3).}
Ligne 15 du bas de la p.~21, remplacer «~inter-section~» par
«~intersection~».

\paragraph{(0, 2.1.6).}
Ligne 13 de la p.~22, remplacer
\[
  Z_i\cap C\left(\bigcup_{j\ne i}Z_j\right)
  \quad\hbox{par}\quad
  \bigcup_{j\ne i}Z_j.
\]

\paragraph{(0, 3.1.6).}
Ligne 17 du bas de la p.~25, remplacer
$\Gamma(U,\sh{F})$ par $\Gamma(V,\sh{F})$.

\paragraph{(0, 4.1.1).}
Ligne 12 de la p.~36, remplacer $\sh{G}\to\sh{F}$ par
$\sh{B}\to\sh{A}$.

\paragraph{(0, 4.2.1).}
Ligne 11 du bas de la p.~39, remplacer «~faisceaux~» par «~faisceau~».

\paragraph{(0, 5.1.2).}
Après la ligne 14 de la p.~45, ajouter :

Si $\sh{G}$ est un $\sh{O}_Y$-Module engendré par ses sections au-dessus
de $Y$, alors $f^*(\sh{G})$ est engendré par ses sections au-dessus de
$X$, car $f^*$ est un foncteur exact à droite.

\paragraph{(0, 5.2.4).}
Ligne 18 du bas de la p.~46, remplacer $f_U$ par $f_U^*$.

\paragraph{(0, 5.3.4).}
Avant la ligne 16 du bas de la p.~47, ajouter :

Si
\[
  \sh{A}\to\sh{B}\to\sh{C}\to\sh{D}\to\sh{E}
\]
est une suite exacte de $\sh{O}_X$-Modules, et si
$\sh{A},\sh{B},\sh{D},\sh{E}$ sont cohérents, alors $\sh{C}$ est
cohérent.

\paragraph{(0, 5.4.1).}
Après la ligne 2 de la p.~49, ajouter :

Supposons $\sh{O}_X$ cohérent, et soit $\sh{F}$ un $\sh{O}_X$-Module
cohérent. Alors, si en un point $x\in X$, $\sh{F}_x$ est un
$\sh{O}_x$-module libre de rang $n$, il existe un voisinage $U$ de $x$
tel que $\sh{F}|U$ soit localement libre de rang $n$; en effet
$\sh{F}_x$ est alors isomorphe à $\sh{O}_x^n$, et la proposition résulte
de (0, 5.2.7).

\paragraph{(0, 6.6.2).}
Remplacer les lignes 10 et 11 de la p.~59 par :

En effet, cela résulte de (0, 6.4.1, d)).

\paragraph{(0, 6.7.1).}
Avant la ligne 10 du bas de la p.~59, ajouter :

Si $f$ est un morphisme plat, on dit encore que $X$ est \emph{plat sur
$Y$}, ou \emph{$Y$-plat}.

\paragraph{(0, 7.2.9).}
Ligne 11 du bas de la p.~65, remplacer $M_n$ par $M_{n-1}$.

\subsection*{Chapitre I}

\paragraph{(I, 1.2.1).}
Ligne 18 de la p.~83, remplacer (deux fois) $K$ par $k$.

\paragraph{(I, 1.2.7).}
Ligne 16 du bas de la p.~84, remplacer $A$ par $A'$ (deux fois) et
$\mathfrak N$ par $\mathfrak N'$; ligne 17 du bas, remplacer $X$ par
$X'$.

\paragraph{(I, 1.7.3) et (I, 2.2.4).}
Les énoncés donnés dans ces numéros peuvent (selon une remarque due à
J.~Tate) être généralisés comme suit :

\subsection*{1.8. Morphismes d'espaces annelés en anneaux locaux dans les schémas affines}
\label{ega2:addendum:subsection:I.1.8-fr}

\begin{proposition}[1.8.1]
\label{ega2:addendum:I.1.8.1-fr}
Soient $(S,\sh{O}_S)$ un schéma affine, $(X,\sh{O}_X)$ un espace annelé
en anneaux locaux. Il y a alors une bijection canonique de l'ensemble
des homomorphismes d'anneaux
\oldpage[II]{218}
\markboth{A. GROTHENDIECK --- Chap. II}{A. GROTHENDIECK --- Chap. II}
\[
  \Gamma(S,\sh{O}_S)\to\Gamma(X,\sh{O}_X)
\]
sur l'ensemble des morphismes d'espaces annelés
$(\psi,\theta):(X,\sh{O}_X)\to(S,\sh{O}_S)$ tels que, pour tout
$x\in X$, $\theta_x^\sharp:\sh{O}_{\psi(x)}\to\sh{O}_x$ soit un
homomorphisme local.
\end{proposition}

Notons d'abord que si $(X,\sh{O}_X)$, $(S,\sh{O}_S)$ sont deux espaces
annelés quelconques, un morphisme $(\psi,\theta)$ de $(X,\sh{O}_X)$
dans $(S,\sh{O}_S)$ définit canoniquement un homomorphisme d'anneaux
$\Gamma(\theta):\Gamma(S,\sh{O}_S)\to\Gamma(X,\sh{O}_X)$; d'où une
première application
\begin{equation}
\label{ega2:addendum:I.1.8.1.1-fr}
  \rho:\Hom((X,\sh{O}_X),(S,\sh{O}_S))
  \to\Hom(\Gamma(S,\sh{O}_S),\Gamma(X,\sh{O}_X)).
\tag{1.8.1.1}
\end{equation}
Inversement, sous les hypothèses de l'énoncé, posons
$A=\Gamma(S,\sh{O}_S)$, et considérons un homomorphisme d'anneaux
$\vphi:A\to\Gamma(X,\sh{O}_X)$. Pour tout $x\in X$, il est clair que
l'ensemble des $f\in A$ tels que $\vphi(f)(x)=0$ est un idéal premier
de $A$, puisque $\sh{O}_x/\mathfrak m_x=\kres(x)$ est un corps; c'est
donc un élément de $S=\Spec(A)$, que nous noterons encore
${}^a\vphi(x)$. En outre, pour tout $f\in A$, on a par définition
(0, 5.5.2)
\[
  {}^a\vphi^{-1}(D(f))=X_f,
\]
ce qui prouve que ${}^a\vphi$ est une application continue $X\to S$.
Définissons ensuite un homomorphisme
\[
  \widetilde{\vphi}:\sh{O}_S\to{}^a\vphi_*(\sh{O}_X)
\]
de $\sh{O}_S$-Modules; pour tout $f\in A$, on a
$\Gamma(D(f),\sh{O}_S)=A_f$ (I, 1.3.6); pour tout $s\in A$, on fera
correspondre à $s/f\in A_f$ l'élément
\[
  (\vphi(s)|X_f)(\vphi(f)|X_f)^{-1}
  \in\Gamma(X_f,\sh{O}_X)
  =\Gamma(D(f),{}^a\vphi_*(\sh{O}_X)),
\]
et on vérifie aussitôt (par passage de $D(f)$ à $D(fg)$) que cela
définit bien un homomorphisme de $\sh{O}_S$-Modules, d'où un morphisme
$({}^a\vphi,\widetilde{\vphi})$ d'espaces annelés. En outre, avec les
mêmes notations, et en posant pour simplifier $y={}^a\vphi(x)$, on voit
aussitôt (0, 3.7.1) que l'on a
\[
  \widetilde{\vphi}_x^\sharp(s_y/f_y)
  =(\vphi(s)_x)(\vphi(f)_x)^{-1};
\]
comme la relation $s_y\in\mathfrak m_y$ est par définition équivalente
à $\vphi(s)_x\in\mathfrak m_x$, on voit que
$\widetilde{\vphi}_x^\sharp$ est un homomorphisme local
$\sh{O}_y\to\sh{O}_x$, et on a ainsi défini une seconde application
\[
  \sigma:\Hom(\Gamma(S,\sh{O}_S),\Gamma(X,\sh{O}_X))\to\mathfrak L,
\]
où $\mathfrak L$ est l'ensemble des morphismes
$(\psi,\theta):(X,\sh{O}_X)\to(S,\sh{O}_S)$ tels que
$\theta_x^\sharp$ soit local pour tout $x\in X$.

Il reste à prouver que $\sigma$ et $\rho$ (restreint à $\mathfrak L$)
sont réciproques l'une de l'autre; or, la définition de
$\widetilde{\vphi}$ montre aussitôt que
$\Gamma(\widetilde{\vphi})=\vphi$, et par suite $\rho\circ\sigma$ est
l'identité. Pour voir que $\sigma\circ\rho$ est l'identité, partons
d'un morphisme $(\psi,\theta)\in\mathfrak L$ et soit
$\vphi=\Gamma(\theta)$; l'hypothèse sur $\theta_x^\sharp$ permet de
déduire de cet homomorphisme, par passage aux quotients, un
monomorphisme
\[
  \theta^x:\kres(\psi(x))\to\kres(x)
\]
tel que pour toute section $f\in A=\Gamma(S,\sh{O}_S)$, on ait
$\theta^x(f(\psi(x)))=\vphi(f)(x)$; la relation $f(\psi(x))=0$ est donc
équivalente à $\vphi(f)(x)=0$, ce qui prouve déjà que
${}^a\vphi=\psi$. D'autre part, les définitions entraînent que le
diagramme
\[
  \xymatrix{
    A\ar[r]^\vphi\ar[d] &
    \Gamma(X,\sh{O}_X)\ar[d]\\
    A_{\psi(x)}\ar[r]^{\theta_x^\sharp} &
    \sh{O}_x
  }
\]
est commutatif, et il en est de même du diagramme analogue où
$\theta_x^\sharp$ est remplacé par
$\widetilde{\vphi}_x^\sharp$, d'où
$\widetilde{\vphi}_x^\sharp=\theta_x^\sharp$ (0, 1.2.4), et par suite
$\widetilde{\vphi}=\theta$.

\begin{env}[1.8.2]
\label{ega2:addendum:I.1.8.2-fr}
Lorsque $(X,\sh{O}_X)$ et $(Y,\sh{O}_Y)$ sont des espaces annelés en
anneaux locaux, nous aurons à considérer les morphismes
$(\psi,\theta):(X,\sh{O}_X)\to(Y,\sh{O}_Y)$ tels que, pour tout
$x\in X$, $\theta_x^\sharp$ soit un homomorphisme local
$\sh{O}_{\psi(x)}\to\sh{O}_x$.

Lorsque nous parlerons désormais
\oldpage[II]{219}
\markboth{ERRATA ET ADDENDA, LISTE 1}{ERRATA ET ADDENDA, LISTE 1}
de \emph{morphisme d'espaces annelés
en anneaux locaux}, c'est toujours d'un tel morphisme qu'il s'agira;
avec cette définition des morphismes, il est clair que les espaces
annelés en anneaux locaux forment une catégorie; pour deux objets
$X,Y$ de cette catégorie, $\Hom(X,Y)$ désignera donc l'ensemble des
morphismes d'espaces annelés en anneaux locaux de $X$ dans $Y$
(ensemble noté $\mathfrak L$ dans (1.8.1)); lorsque nous aurons à
considérer l'ensemble des morphismes d'espaces annelés de $X$ dans
$Y$, nous le noterons $\Hom_{\mathrm{an}}(X,Y)$ pour éviter toute
confusion. L'application (1.8.1.1) s'écrit donc
\begin{equation}
\label{ega2:addendum:I.1.8.2.1-fr}
  \rho:\Hom_{\mathrm{an}}(X,Y)
  \to\Hom(\Gamma(Y,\sh{O}_Y),\Gamma(X,\sh{O}_X)),
\tag{1.8.2.1}
\end{equation}
et sa restriction
\begin{equation}
\label{ega2:addendum:I.1.8.2.2-fr}
  \rho':\Hom(X,Y)
  \to\Hom(\Gamma(Y,\sh{O}_Y),\Gamma(X,\sh{O}_X))
\tag{1.8.2.2}
\end{equation}
est une application fonctorielle en $X$ et $Y$ dans la catégorie des
espaces annelés en anneaux locaux.
\end{env}

\begin{corollary}[1.8.3]
\label{ega2:addendum:I.1.8.3-fr}
Soit $(Y,\sh{O}_Y)$ un espace annelé en anneaux locaux. Pour que $Y$
soit un schéma affine, il faut et il suffit que pour tout espace annelé
en anneaux locaux $(X,\sh{O}_X)$, l'application (1.8.2.2) soit
bijective.
\end{corollary}

La proposition (1.8.1) montre que la condition est nécessaire.
Inversement, si on la suppose vérifiée et si l'on pose
$A=\Gamma(Y,\sh{O}_Y)$, il résulte de l'hypothèse et de (1.8.1) que les
foncteurs $X\mapsto\Hom(X,Y)$ et
$X\mapsto\Hom(X,\Spec(A))$, de la catégorie des espaces annelés en
anneaux locaux dans celle des ensembles, sont isomorphes. On sait que
cela entraîne l'existence d'un isomorphisme canonique
$X\to\Spec(A)$ (cf.~0, 8).

\begin{env}[1.8.4]
\label{ega2:addendum:I.1.8.4-fr}
Soit $S=\Spec(A)$ un schéma affine; désignons par $(S',A')$ l'espace
annelé dont l'espace sous-jacent est réduit à un point et le faisceau
structural $A'$ est le faisceau (nécessairement simple) sur $S'$ défini
par l'anneau $A$. Soit $\pi:S\to S'$ l'unique application de $S$ dans
$S'$; notons d'autre part que pour tout ouvert $U$ de $S$, on a une
application canonique
$\Gamma(S',A')=\Gamma(S,\sh{O}_S)\to\Gamma(U,\sh{O}_S)$ qui définit
donc un $\pi$-morphisme $\iota:A'\to\sh{O}_S$ de faisceaux d'anneaux.
On a ainsi défini canoniquement un morphisme d'espaces annelés
$i=(\pi,\iota):(S,\sh{O}_S)\to(S',A')$. Pour tout $A$-module $M$, nous
désignerons par $M'$ le faisceau simple sur $S'$ défini par $M$, qui
est évidemment un $A'$-Module. Il est clair que l'on a
$i_*(\widetilde M)=M'$ (I, 1.3.7).
\end{env}

\begin{lemma}[1.8.5]
\label{ega2:addendum:I.1.8.5-fr}
Avec les notations de (1.8.4), pour tout $A$-module $M$, le
$\sh{O}_S$-homomorphisme canonique fonctoriel (0, 4.4.3.3)
\begin{equation}
\label{ega2:addendum:I.1.8.5.1-fr}
  i^*(i_*(\widetilde M))\to\widetilde M
\tag{1.8.5.1}
\end{equation}
est un isomorphisme.
\end{lemma}

En effet, les deux membres de (1.8.5.1) sont exacts à droite (le
foncteur $M\mapsto i_*(\widetilde M)$ étant évidemment exact) et
commutent aux sommes directes; en considérant $M$ comme conoyau d'un
homomorphisme $A^{(I)}\to A^{(J)}$, on se ramène à prouver le lemme
dans le cas $M=A$, et il est évident dans ce cas.

\begin{corollary}[1.8.6]
\label{ega2:addendum:I.1.8.6-fr}
Soient $(X,\sh{O}_X)$ un espace annelé, $u:X\to S$ un morphisme
d'espaces
\oldpage[II]{220}
\markboth{A. GROTHENDIECK --- Chap. II}{A. GROTHENDIECK --- Chap. II}
annelés. Pour tout $A$-module $M$, on a (avec les notations de (1.8.4)) un
isomorphisme canonique fonctoriel de $\sh{O}_X$-Modules
\begin{equation}
\label{ega2:addendum:I.1.8.6.1-fr}
  u^*(\widetilde M)\isoto u^*(i^*(M')).
\tag{1.8.6.1}
\end{equation}
\end{corollary}

\begin{corollary}[1.8.7]
\label{ega2:addendum:I.1.8.7-fr}
Sous les hypothèses de (1.8.6), on a, pour tout $A$-module $M$ et tout
$\sh{O}_X$-Module $\sh{F}$, un isomorphisme canonique fonctoriel en
$M$ et $\sh{F}$
\begin{equation}
\label{ega2:addendum:I.1.8.7.1-fr}
  \Hom_{\sh{O}_S}(\widetilde M,u_*(\sh{F}))
  \isoto\Hom_A(M,\Gamma(X,\sh{F})).
\tag{1.8.7.1}
\end{equation}
\end{corollary}

On a en effet, en vertu de (0, 4.4.3) et du lemme (1.8.5), un
isomorphisme canonique de bifoncteurs
\[
  \Hom_{\sh{O}_S}(\widetilde M,u_*(\sh{F}))
  \isoto\Hom_{A'}(M',i_*(u_*(\sh{F}))),
\]
et il est clair que le second membre n'est autre que
$\Hom_A(M,\Gamma(X,\sh{F}))$. On notera que l'homomorphisme canonique
(1.8.7.1) fait correspondre à tout $\sh{O}_S$-homomorphisme
$h:\widetilde M\to u_*(\sh{F})$ (autrement dit, à tout $u$-morphisme
$\widetilde M\to\sh{F}$) le $A$-homomorphisme
\[
  \Gamma(h):M\to\Gamma(S,u_*(\sh{F}))=\Gamma(X,\sh{F}).
\]

\begin{env}[1.8.8]
\label{ega2:addendum:I.1.8.8-fr}
Avec les notations de (1.8.4), il est clair (0, 4.1.1) que tout
morphisme d'espaces annelés $(\psi,\theta):X\to S'$ équivaut à la
donnée d'un homomorphisme d'anneaux
$A\to\Gamma(X,\sh{O}_X)$. On peut donc encore interpréter (1.8.1)
comme définissant une bijection canonique
\[
  \Hom(X,S)\isoto\Hom(X,S')
\]
(où bien entendu il s'agit au second membre de morphismes d'espaces
annelés, puisqu'en général $A$ n'est pas un anneau local). Plus
généralement, si $X,Y$ sont deux espaces annelés en anneaux locaux et
si $(Y',A')$ est l'espace annelé dont l'espace sous-jacent est réduit à
un point et dont le faisceau d'anneaux $A'$ est le faisceau simple
défini par l'anneau $\Gamma(Y,\sh{O}_Y)$, on peut interpréter
(1.8.2.1) comme une application
\begin{equation}
\label{ega2:addendum:I.1.8.8.1-fr}
  \rho:\Hom_{\mathrm{an}}(X,Y)\to\Hom(X,Y').
\tag{1.8.8.1}
\end{equation}
Le résultat de (1.8.3) s'interprète donc en disant que les schémas
affines sont caractérisés parmi les espaces annelés en anneaux locaux
comme ceux pour lesquels la restriction de $\rho$ à $\Hom(X,Y)$ :
\begin{equation}
\label{ega2:addendum:I.1.8.8.2-fr}
  \rho':\Hom(X,Y)\to\Hom(X,Y')
\tag{1.8.8.2}
\end{equation}
est bijective pour tout espace annelé en anneaux locaux $X$. Dans un
chapitre ultérieur, nous généraliserons cette définition, ce qui
permettra d'associer à un espace annelé quelconque $Z$ (et non plus
seulement à un espace annelé dont l'espace sous-jacent est réduit à un
point) un espace annelé en anneaux locaux que nous appellerons encore
$\Spec(Z)$; cela sera le point de départ d'une théorie «~relative~» des
préschémas au-dessus d'espaces annelés quelconques, étendant les
résultats du chap.~I.
\end{env}

\begin{env}[1.8.9]
\label{ega2:addendum:I.1.8.9-fr}
On peut considérer les couples $(X,\sh{F})$ formés d'un espace annelé
en anneaux locaux $X$ et d'un $\sh{O}_X$-Module $\sh{F}$ comme formant
une catégorie, un morphisme $(X,\sh{F})\to(Y,\sh{G})$ de cette catégorie
étant un couple $(u,h)$ formé d'un morphisme d'espaces
\oldpage[II]{221}
\markboth{ERRATA ET ADDENDA, LISTE 1}{ERRATA ET ADDENDA, LISTE 1}
annelés en anneaux locaux $u:X\to Y$ et d'un $u$-morphisme
$h:\sh{G}\to\sh{F}$ de Modules; ces
morphismes (pour $(X,\sh{F})$ et $(Y,\sh{G})$ fixés) forment un ensemble
que nous noterons $\Hom((X,\sh{F}),(Y,\sh{G}))$; l'application
$(u,h)\mapsto(\rho'(u),\Gamma(h))$ est une application canonique
\begin{equation}
\label{ega2:addendum:I.1.8.9.1-fr}
\begin{split}
  \Hom((X,\sh{F}),(Y,\sh{G}))\to
  \Hom\bigl(
    &(\Gamma(Y,\sh{O}_Y),\Gamma(Y,\sh{G})),\\
    &(\Gamma(X,\sh{O}_X),\Gamma(X,\sh{F}))
  \bigr),
\end{split}
\tag{1.8.9.1}
\end{equation}
fonctorielle en $(X,\sh{F})$ et $(Y,\sh{G})$, le second membre étant
l'ensemble des di-homomorphismes correspondant aux anneaux et aux
modules considérés (0, 1.0.2).
\end{env}

\begin{corollary}[1.8.10]
\label{ega2:addendum:I.1.8.10-fr}
Soient $Y$ un espace annelé en anneaux locaux, $\sh{G}$ un
$\sh{O}_Y$-Module. Pour que $Y$ soit un schéma affine et $\sh{G}$ un
$\sh{O}_Y$-Module quasi-cohérent, il faut et il suffit que pour tout
couple $(X,\sh{F})$ formé d'un espace annelé en anneaux locaux $X$ et
d'un $\sh{O}_X$-Module $\sh{F}$, l'application canonique (1.8.8.1) soit
bijective.
\end{corollary}

Nous laissons au lecteur le soin de développer le raisonnement, calqué
sur celui de (1.8.3), et où on utilise (1.8.1) et (1.8.7).

\begin{remark}[1.8.11]
\label{ega2:addendum:I.1.8.11-fr}
Les énoncés (I, 1.7.3), (I, 1.7.4) et (I, 2.2.4) sont des cas
particuliers de (1.8.1), ainsi que la définition (I, 1.6.1); de même
(I, 2.2.5) résulte de (1.8.7). La proposition (1.8.7) entraîne aussi
(I, 1.6.3) (et par suite (I, 1.6.4)) comme cas particulier, car si $X$
est affine et $\Gamma(X,\sh{F})=N$, les foncteurs
\[
  M\longmapsto
  \Hom_{\sh{O}_S}(\widetilde M,u_*(\widetilde N))
  \quad\hbox{et}\quad
  M\longmapsto
  \Hom_{\sh{O}_S}(\widetilde M,(N_{[\vphi]})\supertilde)
\]
(où $\vphi:A\to\Gamma(X,\sh{O}_X)$ correspond à $u$) sont isomorphes
par (1.8.7) et (I, 1.3.8). Enfin, (I, 1.6.5) (et par suite (I, 1.6.6))
résulte de (1.8.6) et du fait que pour tout $f\in A$, les $A_f$-modules
\[
  N'\otimes_{A'}A_f
  \quad\hbox{et}\quad
  (N'\otimes_{A'}A)_f
\]
(notations de (I, 1.6.5)) sont canoniquement isomorphes.
\end{remark}

\paragraph{Insérer après (I, 3.2.8) :}

\begin{env}[3.2.9]
\label{ega2:addendum:I.3.2.9-fr}
Il résulte de (1.8.1) que l'on peut préciser (I, 3.2.2) de la façon
suivante : $Z=\Spec(B\otimes_A C)$ est non seulement un produit de
$X=\Spec(B)$ et de $Y=\Spec(C)$ dans la catégorie des $S$-préschémas,
mais aussi dans la catégorie des espaces annelés en anneaux locaux,
au-dessus de $S$ (avec une définition des $S$-morphismes calquée sur
celle de (I, 2.5.2)). La démonstration de (I, 3.2.6) prouve alors en
fait que pour deux $S$-préschémas quelconques $X,Y$, le préschéma
$X\times_S Y$ est non seulement le produit de $X$ et $Y$ dans la
catégorie des $S$-préschémas, mais encore dans la catégorie des espaces
annelés en anneaux locaux au-dessus du préschéma $S$.
\end{env}

\paragraph{(I, 4.4.5).}
Ligne 14 du bas de la p.~126, remplacer $B$ par $A$; ligne 13 du bas,
remplacer «~$A$-algèbre~» par «~$B$-algèbre~».

\paragraph{(I, 5.3.5).}
Ligne 2 du bas de la p.~132, remplacer $Y\to S$ par $g:Y\to S$.

\paragraph{(I, 6.3.2.1).}
Ligne 6 du bas de la p.~144, lire :
\[
  D(g_i)\subset W.
\]

\paragraph{(I, 7.3.8).}
Remplacer le texte actuel, qui est erroné, par le suivant :

Soient $X,Y$ deux préschémas intègres, ce qui entraîne que
$\sh{R}(X)$ (resp. $\sh{R}(Y)$) est un $\sh{O}_X$-Module (resp. un
$\sh{O}_Y$-Module) quasi-cohérent (I, 7.3.3). Soit $f:X\to Y$ un
morphisme dominant; alors il existe un homomorphisme canonique de
$\sh{O}_X$-Modules
\oldpage[II]{222}
\markboth{A. GROTHENDIECK --- Chap. II}{A. GROTHENDIECK --- Chap. II}
\begin{equation}
\label{ega2:addendum:I.7.3.8.1-fr}
  \tau:f^*(\sh{R}(Y))\to\sh{R}(X).
\tag{7.3.8.1}
\end{equation}

Supposons d'abord $X=\Spec(A)$ et $Y=\Spec(B)$ affines d'anneaux
intègres $A$ et $B$, $f$ correspondant donc à un homomorphisme injectif
$B\to A$, qui se prolonge en un monomorphisme $L\to K$ du corps des
fractions $L$ de $B$ dans le corps des fractions $K$ de $A$.
L'homomorphisme (7.3.8.1) correspond alors à l'homomorphisme canonique
\[
  L\otimes_B A\to K
\]
(I, 1.6.5). Dans le cas général, pour tout couple d'ouverts affines non
vides $U\subset X$, $V\subset Y$ tels que $f(U)\subset V$, on définit
de la façon précédente un homomorphisme $\tau_{U,V}$ et on constate
aussitôt que si $U'\subset U$, $V'\subset V$,
$f(U')\subset V'$, $\tau_{U,V}$ prolonge $\tau_{U',V'}$, d'où notre
assertion. Si $x,y$ sont les points génériques de $X$ et $Y$
respectivement, on a $f(x)=y$,
\[
  \bigl(f^*(\sh{R}(Y))\bigr)_x
  =\sh{O}_y\otimes_{\sh{O}_y}\sh{O}_x
  =\sh{O}_x
\]
(0, 4.3.1), et $\tau_x$ est donc un isomorphisme.

\paragraph{(I, 9.5.2).}
Dans les 3 dernières lignes de la p.~176 et les 5 premières de la
p.~177, remplacer partout $X'$ par $Y'$ et $X''$ par $Y''$.

\paragraph{(I, 10.5.4).}
Ligne 11 du bas de la p.~187, remplacer (10.3.4) par (10.3.5).

\paragraph{(I, 10.6.3).}
Ligne 3 du bas de la p.~189, remplacer $X$ par $\mathfrak X$.

\paragraph{(I, 10.14.2).}
Ligne 5 de la p.~210, remplacer «~$\sh{O}_{\mathfrak X}$-Module cohérent~» par
«~Idéal cohérent de $\sh{O}_{\mathfrak X}$~»; ligne 7, remplacer
«~$\sh{O}_{\mathfrak X}$-Modules cohérents~» par «~Idéaux cohérents de
$\sh{O}_{\mathfrak X}$~».
